\documentclass[pdflatex,iicol,sn-basic]{sn-jnl}

\usepackage{amsmath,amssymb,mathtools}
\usepackage{booktabs}
\usepackage{microtype}
\usepackage{enumitem}
\usepackage[nameinlink,noabbrev]{cleveref}
\usepackage{url}

\hypersetup{
  colorlinks=true,
  linkcolor=blue,
  citecolor=blue,
  urlcolor=blue,
  pdftitle={Sequential Depletion Ordering with Residual-Fraction Costs},
  pdfauthor={Joseph L. Malone}
}

\theoremstyle{thmstyleone}
\newtheorem{theorem}{Theorem}[section]
\newtheorem{proposition}[theorem]{Proposition}
\newtheorem{lemma}[theorem]{Lemma}
\newtheorem{corollary}[theorem]{Corollary}

\theoremstyle{thmstyletwo}
\newtheorem{example}[theorem]{Example}

\theoremstyle{thmstylethree}
\newtheorem{remark}[theorem]{Remark}
\newtheorem{definition}[theorem]{Definition}

\newcommand{\R}{\mathbb{R}}
\newcommand{\E}{\mathbb{E}}
\newcommand{\cI}{\mathcal{I}}
\newcommand{\lr}{\le_{\mathrm{lr}}}

\begin{document}

\title[Sequential Depletion Ordering]{Sequential Depletion Ordering with
Residual-Fraction Costs}

\author{\fnm{Joseph L.} \sur{Malone}\par
\vspace{2pt}
{\normalfont\small Contact: m[A]l[o]n\{e\}(a)salient.vision}}

\abstract{
A fixed collection of positive loads is removed sequentially from a resource
with permanent reserve \(q\ge 0\). If removing size \(x\) at residual level
\(R\) costs the residual fraction \(x/R\), the cost of a permutation \(\pi\) is
\[
C_q(\pi)
=
\sum_{k=1}^{n}
\frac{x_{\pi(k)}}{q+\sum_{j=k}^{n}x_{\pi(j)}}.
\]
An exact adjacent-interchange identity shows that nondecreasing size order
minimizes \(C_q\) and nonincreasing order maximizes it. For \(q>0\), the
log-depletion vector \(d_k=\log(R_k/R_{k+1})\) is majorized between the
decreasing-size and increasing-size extremes, so every Schur-convex functional
of \(d\) shares those extremes. The same comparison classifies nonlinear stage
costs \(\sum\phi(x_{\pi(k)}/R_k)\) by the curvature of
\(g(t)=\phi(1-e^{-t})\).

With precedence constraints, an order-ideal dynamic program computes the
optimum and smallest-available greedy can fail. For independent sizes totally
ordered by monotone likelihood ratio, the ex-ante small-first schedule
minimizes expected cost. For multi-resource loads, componentwise chains retain
the sorting rule, while heterogeneous vectors admit residual-state preference
reversals. Related residual-fraction sums appear in cumulative Riemann-sum
bounds; those results give order-independent integral estimates rather than a
permutation theory. The contribution is the ordering, majorization, constrained,
stochastic, and network analysis built on one interchange identity.
}

\keywords{Scheduling, residual fraction, adjacent interchange, majorization,
stochastic orders, precedence constraints}

\maketitle

\section{Introduction}
\label{sec:intro}

Suppose a system has residual size \(R\) and a load of size \(x\) is removed.
The relative shock is \(x/R\). If several fixed loads must be removed in some
order, each removal changes the residual, so the sequence of shocks depends on
the order.

Let \(x_1,\dots,x_n>0\) and let \(q\ge 0\) be reserve that is never removed.
For a permutation \(\pi\) of \([n]=\{1,\ldots,n\}\), define
\begin{equation}\label{eq:objective}
C_q(\pi)
=
\sum_{k=1}^{n}
\frac{x_{\pi(k)}}{q+\sum_{j=k}^{n}x_{\pi(j)}}.
\end{equation}
General sum-of-ratios programs are hard \citep{SchaibleShi2003}. The suffix
structure in \eqref{eq:objective} is special: an exact adjacent-interchange
identity determines the optimum.

The main facts are these.
\begin{enumerate}[leftmargin=2em]
  \item Nondecreasing size order minimizes \(C_q\); nonincreasing order maximizes it.
  \item For \(q>0\), the log-depletion vectors of all permutations lie between
  the decreasing-size and increasing-size extremes in the majorization order.
  \item Stage costs \(\phi(x/R)\) inherit those extremal orders according to the
  curvature of \(g(t)=\phi(1-e^{-t})\).
  \item With precedence constraints, the optimum is given by a dynamic program
  over order ideals; the natural smallest-available rule is not always optimal.
  \item If independent random sizes are totally ordered by monotone likelihood
  ratio, the corresponding fixed order minimizes expected cost.
  \item Componentwise-comparable multi-resource loads still sort; for
  heterogeneous loads, even the locally preferred order can depend on the
  residual state.
\end{enumerate}

The deterministic swap argument is classical adjacent interchange in the sense
of \citet{Smith1956}, \citet{Sidney1975}, and \citet{Lawler1978}. Closely
related residual-ratio sums also occur as cumulative Riemann sums
\citep{Pain2026}. For \(q>0\), reverse normalization identifies
\eqref{eq:objective} exactly with one such sum and yields an order-independent
integral bound. That result does not compare permutations. The contribution of
this paper is the permutation theory for \eqref{eq:objective} and the
majorization, constrained, stochastic, and multi-resource results built on its
interchange identity.

\section{Related work}
\label{sec:related}

\subsection{Adjacent interchange and scheduling}
Local pairwise comparisons are standard in single-machine sequencing
\citep{Smith1956,Sidney1975,Lawler1978,CorreaSchulz2005,JagerWarode2024}. Those
arguments usually target completion-time objectives. The cost
\eqref{eq:objective} is a nested residual-fraction sum, not weighted completion
time.

\subsection{Cumulative residual fractions}
\citet{Pain2026} studies sums \(\sum_{i=1}^n a_i g(S_i)\) with \(a_i>0\),
\(\sum_i a_i=1\), and prefix totals \(S_i=\sum_{j=1}^i a_j\). For decreasing
integrable \(g\), that work proves
\[
\sum_{i=1}^n a_i g(S_i)\le\int_0^1 g(x)\,dx,
\]
an order-independent integral upper bound. To make the connection exact, let
\(X=\sum_i x_i\), set \(\rho=q/X\), reverse the schedule via
\(\sigma(i)=\pi(n+1-i)\), and define
\[
a_i=\frac{x_{\sigma(i)}}{X},
\qquad
S_i=\sum_{j=1}^i a_j.
\]
Then, for \(q>0\),
\begin{equation}\label{eq:pain}
C_q(\pi)
=
\sum_{i=1}^n \frac{a_i}{\rho+S_i}
\le
\int_0^1\frac{ds}{\rho+s}
=
\log\left(1+\frac{X}{q}\right).
\end{equation}
The inequality is Pain's theorem with \(g(s)=1/(\rho+s)\). It controls every
order but neither compares two orders nor identifies an extremal permutation.
The case \(q=0\) has the same cumulative form with \(g(s)=1/s\), but that
function is not integrable at zero and is outside the stated integral-bound
theorem.

\subsection{Majorization}
Majorization and Karamata's inequality are classical
\citep{HardyLittlewoodPolya1952,MarshallOlkinArnold2011,Arnold2008}. Here they
apply to the vector of log residual ratios generated by a permutation.

\subsection{Sequential choice and size-biased order}
Stagewise terms \(x_i/\sum_{j\in S}x_j\) appear as choice probabilities in Luce
and Plackett--Luce models and in size-biased residual allocation
\citep{Luce1959,Plackett1975,PermanPitmanYor1992,PitmanTran2015}. That work
studies random orders and products of choice probabilities rather than
optimization of the deterministic sum \eqref{eq:objective}.

\subsection{Stochastic orders and rearrangement}
Monotone likelihood ratio is a standard stochastic order
\citep{KarlinRubin1956,ShakedShanthikumar2007}. Bivariate characterizations and
stochastic rearrangement inequalities provide general foundations for pairwise
interchange under likelihood-ratio order
\citep{ShanthikumarYao1991,ChangYao1993}. This paper applies that established
method to the explicit swap kernel of \(C_q\).

\subsection{Precedence and order ideals}
Dynamic programming over order ideals is a standard exact method for separable
costs on linear extensions of a poset \citep{Steiner1984}. The recursion below
specializes that classical state space to residual-fraction stage costs, with a
width-based complexity bound and a greedy counterexample for this objective.

\section{The deterministic depletion objective}
\label{sec:deterministic}

For a permutation \(\pi\), write
\[
R_k(\pi)=q+\sum_{j=k}^{n}x_{\pi(j)}
\]
for the resource present immediately before stage \(k\). Because every size is
positive, \(R_k(\pi)\ge x_{\pi(k)}>0\) for all \(k\), including the case
\(q=0\).

\begin{lemma}[Exact adjacent-interchange identity]
\label{lem:swap}
Suppose two adjacent items have sizes \(a,b>0\), and let \(T\ge 0\) be the
reserve plus the total size of all items scheduled after them. Then
\begin{align}
&C_q(\ldots,a,b,\ldots)-C_q(\ldots,b,a,\ldots)\nonumber\\
&\qquad=
\frac{ab(a-b)}{(T+a+b)(T+a)(T+b)}.
\label{eq:swap}
\end{align}
\end{lemma}

\begin{proof}
Only the two displayed stages change. Their difference is
\begin{align*}
&\left(\frac{a}{T+a+b}+\frac{b}{T+b}\right)
-
\left(\frac{b}{T+a+b}+\frac{a}{T+a}\right)\\
&=\frac{a-b}{T+a+b}
+\frac{T(b-a)}{(T+a)(T+b)}\\
&=\frac{ab(a-b)}{(T+a+b)(T+a)(T+b)}.
\end{align*}
All denominators are positive because \(a,b>0\) and \(T\ge 0\).
\end{proof}

\begin{theorem}[Extremal orders]
\label{thm:sort}
For every \(q\ge 0\), nondecreasing size order minimizes \(C_q\), and
nonincreasing size order maximizes \(C_q\). The minimizer and maximizer are
unique up to permutations of equal-sized items.
\end{theorem}

\begin{proof}
If an adjacent pair is inverted, \(a>b\), then \cref{lem:swap} shows that
swapping it from \((a,b)\) to \((b,a)\) strictly decreases the objective.
Repeatedly removing inversions yields nondecreasing order. Reversing the
argument yields nonincreasing order as the maximizer.
\end{proof}

When \(q=0\), the same theorem gives a harmonic bound.

\begin{corollary}[Harmonic sandwich]
\label{cor:harmonic}
Let \(H_n=\sum_{k=1}^n 1/k\). For \(q=0\),
\[
C_0(x^{\uparrow})\le H_n\le C_0(x^{\downarrow}),
\]
with equality on either side if and only if all sizes are equal.
\end{corollary}

\begin{proof}
For \(y_1\le\cdots\le y_n\),
\[
\sum_{j=k}^n y_j\ge(n-k+1)y_k,
\]
so the \(k\)th term is at most \(1/(n-k+1)\). The descending inequality is
analogous.
\end{proof}

\section{Log-depletion majorization}
\label{sec:majorization}

Assume throughout this section that \(q>0\), so every residual is strictly
larger than the load about to be removed and every log increment is finite.
Define the fractional shock and log-depletion shock at stage \(k\) by
\begin{equation}\label{eq:z-d}
\begin{aligned}
z_k(\pi)&=\frac{x_{\pi(k)}}{R_k(\pi)},\\
d_k(\pi)&=\log\frac{R_k(\pi)}{R_{k+1}(\pi)}
=-\log(1-z_k(\pi)).
\end{aligned}
\end{equation}
The total log depletion telescopes:
\begin{equation}\label{eq:fixedsum}
\sum_{k=1}^n d_k(\pi)
=\log\frac{q+\sum_i x_i}{q},
\end{equation}
independently of \(\pi\).

For vectors \(u,v\in\R^n\) with equal coordinate sums, write \(u\prec v\) when
\(u\) is majorized by \(v\); equivalently, the sums of the \(m\) largest
coordinates of \(u\) do not exceed those of \(v\) for \(m<n\)
\citep{MarshallOlkinArnold2011}.

\begin{theorem}[Log-depletion majorization sandwich]
\label{thm:maj}
Let \(d(\pi)=(d_1(\pi),\ldots,d_n(\pi))\). Then every permutation satisfies
\begin{equation}\label{eq:maj}
d(x^{\downarrow})\prec d(\pi)\prec d(x^{\uparrow}).
\end{equation}
Thus decreasing size order produces the least concentrated log-depletion
vector, and increasing size order produces the most concentrated one.
\end{theorem}

\begin{proof}
Consider an adjacent inversion \(a\ge b\) with tail resource \(T>0\). Before
swapping, the two log shocks are
\[
u_1=\log\frac{T+a+b}{T+b},
\qquad
u_2=\log\frac{T+b}{T}.
\]
After swapping to \((b,a)\), they are
\[
v_1=\log\frac{T+a+b}{T+a},
\qquad
v_2=\log\frac{T+a}{T}.
\]
The pair sums are equal. Also
\[
\begin{aligned}
u_2&\le v_2,\\
u_1&=\log\left(1+\frac{a}{T+b}\right)\\
&\le\log\left(1+\frac{a}{T}\right)=v_2.
\end{aligned}
\]
Hence \(\max\{u_1,u_2\}\le\max\{v_1,v_2\}\). For two-dimensional vectors of
equal sum, this is exactly \((u_1,u_2)\prec(v_1,v_2)\). Replacing two
coordinates by a vector that majorizes them preserves majorization after
adjoining all unchanged coordinates. Repeated swaps from decreasing to
increasing size order prove \eqref{eq:maj}.
\end{proof}

\begin{corollary}[Universal convex-risk ordering]
\label{cor:convex}
For every convex function \(g\) on \([0,\infty)\),
\[
\sum_k g(d_k(x^{\downarrow}))
\le
\sum_k g(d_k(\pi))
\le
\sum_k g(d_k(x^{\uparrow})).
\]
For concave \(g\), both inequalities reverse.
\end{corollary}

\begin{proof}
Apply Karamata's inequality to \cref{thm:maj}
\citep{HardyLittlewoodPolya1952,MarshallOlkinArnold2011}.
\end{proof}

Consequently, decreasing size order minimizes the maximum log shock, every
\(\ell_p\) norm with \(p>1\), the variance of the log-shock vector, and every
separable convex risk measure on \(d\). Increasing order minimizes the raw
fraction sum because \(z=1-e^{-d}\) is concave in \(d\).

\subsection{Nonlinear stage penalties and a phase transition}

Let \(\phi:[0,1)\to\R\) and define
\begin{equation}\label{eq:nonlinear}
J_{\phi,q}(\pi)=\sum_{k=1}^n\phi(z_k(\pi)).
\end{equation}
Set
\[
g_\phi(t)=\phi(1-e^{-t}).
\]

\begin{corollary}[Curvature classification]
\label{cor:curvature}
If \(g_\phi\) is concave, increasing size order minimizes \(J_{\phi,q}\) and
decreasing size order maximizes it. If \(g_\phi\) is convex, decreasing size
order minimizes \(J_{\phi,q}\) and increasing size order maximizes it.
\end{corollary}

If \(\phi\) is twice differentiable and \(z=1-e^{-t}\), then
\begin{equation}\label{eq:curvature}
g_\phi''(t)
=(1-z)\bigl((1-z)\phi''(z)-\phi'(z)\bigr).
\end{equation}
The order-invariant boundary solves
\[
(1-z)\phi''(z)=\phi'(z),
\]
so
\begin{equation}\label{eq:boundary}
\phi(z)=A[-\log(1-z)]+B.
\end{equation}

An explicit family makes the transition concrete. For \(\alpha\in\R\), define
\begin{equation}\label{eq:alpha}
\phi_\alpha(z)=
\begin{cases}
\displaystyle\frac{(1-z)^{-\alpha}-1}{\alpha},&\alpha\ne0,\\[6pt]
-\log(1-z),&\alpha=0.
\end{cases}
\end{equation}
Then
\[
\phi_\alpha(1-e^{-t})=
\begin{cases}
(e^{\alpha t}-1)/\alpha,&\alpha\ne0,\\
t,&\alpha=0,
\end{cases}
\]
and its second derivative has the sign of \(\alpha\).

\begin{corollary}[Small-first / log-invariant / large-first transition]
\label{cor:phase}
For the family \eqref{eq:alpha}:
\begin{enumerate}[label=(\roman*),leftmargin=2em]
\item if \(\alpha<0\), increasing size order minimizes \(J_{\phi_\alpha,q}\);
\item if \(\alpha=0\), every order has identical cost;
\item if \(\alpha>0\), decreasing size order minimizes \(J_{\phi_\alpha,q}\).
\end{enumerate}
In particular, \(\alpha=-1\) recovers the raw fraction \(\phi(z)=z\), while
\(\alpha=1\) recovers the post-removal odds \(\phi(z)=z/(1-z)\).
\end{corollary}

Three natural ways to score depletion therefore sit on opposite sides of the
same line: residual share before removal, pure log change, and residual share
after removal.

\section{Precedence-constrained depletion}
\label{sec:precedence}

Let \(P=(V,\preceq)\) be a finite precedence poset. A feasible schedule is a
linear extension of \(P\). For \(S\subseteq V\), let
\[
Q(S)=q+\sum_{j\in V\setminus S}x_j.
\]
A set \(S\) is an order ideal if \(j\in S\) and \(i\preceq j\) imply \(i\in S\).
Let \(A(S)\) be the set of minimal elements of \(V\setminus S\); these are the
jobs available next. As in the unconstrained case, every size is positive, so
\(Q(S)>0\) for every proper ideal \(S\).

\begin{theorem}[Exact dynamic program over ideals]
\label{thm:dp}
Define \(F(V)=0\) and, for every proper order ideal \(S\),
\begin{equation}\label{eq:dp}
F(S)=\min_{i\in A(S)}
\left\{
\frac{x_i}{Q(S)}+F(S\cup\{i\})
\right\}.
\end{equation}
Then \(F(\varnothing)\) is the minimum of \(C_q\) over all linear extensions of
\(P\).
\end{theorem}

\begin{proof}
Every feasible continuation from \(S\) must choose an available minimal element
\(i\in A(S)\). Its immediate denominator is \(Q(S)\), independent of which
available job is selected. After choosing \(i\), the processed set is the ideal
\(S\cup\{i\}\). The principle of optimality gives \eqref{eq:dp}.
\end{proof}

Let \(\cI(P)\) be the set of order ideals and let \(w\) be the width of \(P\).

\begin{proposition}[Width-sensitive complexity]
\label{prop:width}
With available sets maintained incrementally, \eqref{eq:dp} can be evaluated in
\(O(w|\cI(P)|)\) state transitions, up to input-preprocessing costs. If a
width-\(w\) chain decomposition has chain lengths \(n_1,\ldots,n_w\), then
\begin{equation}\label{eq:idealsbound}
|\cI(P)|\le\prod_{r=1}^w(n_r+1),
\end{equation}
so the algorithm is polynomial for every fixed width \(w\).
\end{proposition}

\begin{proof}
The available set is an antichain, hence has size at most \(w\). An ideal
intersects each chain in a prefix, and therefore injects into the tuple of its
\(w\) prefix lengths, proving \eqref{eq:idealsbound}.
\end{proof}

In general this is exponential in \(n\). For fixed width it is \(n^{O(w)}\). The
unconstrained increasing-order cost of the remaining jobs is also an admissible
lower bound for branch-and-bound.

\begin{lemma}[Local inversion exclusion]
\label{lem:local}
No optimal linear extension contains adjacent incomparable jobs \(i,j\) with
\(x_i>x_j\) in that order.
\end{lemma}

\begin{proof}
Because the jobs are incomparable and adjacent, swapping them preserves
feasibility. The swap strictly reduces cost by \cref{lem:swap}.
\end{proof}

Local inversion exclusion does not justify a smallest-available greedy policy.

\begin{example}[Smallest available can be suboptimal]
\label{ex:greedy}
Let \(q=1\) and let \(A\prec C\), with \(B\) incomparable to both. Set
\[
(x_A,x_B,x_C)=(8,7,2).
\]
Initially, \(A\) and \(B\) are available, so smallest-available chooses \(B\),
producing
\[
C_1(B,A,C)=\frac7{18}+\frac8{11}+\frac23\approx1.7828.
\]
The feasible order \((A,C,B)\) has
\[
C_1(A,C,B)=\frac8{18}+\frac2{10}+\frac78\approx1.5194.
\]
Processing the larger predecessor first unlocks the very small successor and is
globally preferable.
\end{example}

\section{Stochastic sizes under likelihood-ratio order}
\label{sec:stochastic}

Suppose sizes are independent positive random variables \(X_1,\ldots,X_n\), and
the schedule must be selected ex ante. Let all variables admit densities with
respect to a common dominating measure. Recall that \(X\lr Y\) if
\(f_Y(x)/f_X(x)\) is nondecreasing wherever defined
\citep{ShakedShanthikumar2007}.

\begin{theorem}[Stochastic small-first rule under MLR order]
\label{thm:stochastic}
Assume
\[
X_1\lr X_2\lr\cdots\lr X_n.
\]
Then the order \((1,2,\ldots,n)\) minimizes \(\E[C_q(\pi)]\) over all fixed
permutations \(\pi\). Reversing the order maximizes the expectation. Under
strict likelihood-ratio inequalities on sets of positive measure, the extremal
orders are unique up to identically distributed jobs.
\end{theorem}

\begin{proof}
It suffices to compare an adjacent pair \(X,Y\) with \(X\lr Y\). Let \(T\ge0\)
be the reserve plus the sum of all random sizes after the pair. Independence
makes \(T\) independent of \((X,Y)\). Conditional on \(T=t\), \cref{lem:swap}
gives
\begin{equation}\label{eq:hkernel}
h_t(x,y)
=
\frac{xy(x-y)}{(t+x+y)(t+x)(t+y)}
\end{equation}
after subtracting the cost of order \((Y,X)\) from that of order \((X,Y)\). The
kernel is antisymmetric and positive when \(x>y\). Pairing the regions \(x>y\)
and \(y>x\) yields
\begin{align*}
\E[h_t(X,Y)]
&=\int_{x>y} h_t(x,y)\\
&\quad{}\times
\bigl[f_X(x)f_Y(y)-f_X(y)f_Y(x)\bigr]\,dx\,dy.
\end{align*}
For \(x>y\), likelihood-ratio order implies
\[
\frac{f_Y(x)}{f_X(x)}\ge\frac{f_Y(y)}{f_X(y)},
\]
so the bracketed factor is nonpositive. Therefore \(\E[h_t(X,Y)]\le0\) for every
\(t\), and hence
\[
\E[C_q(\ldots,X,Y,\ldots)]
\le
\E[C_q(\ldots,Y,X,\ldots)].
\]
Adjacent interchanges sort the jobs into likelihood-ratio order. Reversing each
comparison gives the maximizer.
\end{proof}

\begin{corollary}[Common parametric families]
\label{cor:families}
The stochastic theorem applies to, among others:
\begin{enumerate}[label=(\roman*),leftmargin=2em]
\item exponential sizes, ordered by decreasing rate;
\item gamma sizes with common shape, ordered by decreasing rate;
\item lognormal sizes with common variance, ordered by increasing log-location
parameter.
\end{enumerate}
\end{corollary}

Ordering by expected size is not enough. Likelihood-ratio order controls the
sign of the full antisymmetric swap kernel in every tail state.

\section{Multi-resource and network depletion}
\label{sec:network}

Let there be \(m\) resource pools. Item \(i\) has a nonnegative load vector
\[
x_i=(x_{i1},\ldots,x_{im})\in\R_+^m
\]
that is not the zero vector, resource \(r\) has reserve \(q_r>0\), and
\(w_r>0\) is its importance weight. The network objective is
\begin{equation}\label{eq:networkobjective}
C^{\mathrm{net}}(\pi)
=
\sum_{r=1}^m w_r
\sum_{k=1}^{n}
\frac{x_{\pi(k),r}}
{q_r+\sum_{j=k}^n x_{\pi(j),r}},
\end{equation}
The positive reserves keep every resourcewise denominator well defined even
when an item has zero load on a particular resource.

\begin{lemma}[Network swap identity]
\label{lem:netswap}
For adjacent load vectors \(a,b\in\R_+^m\) and tail resources \(T_r>0\),
\begin{equation}\label{eq:netswap}
\Delta(a,b;T)
=
\sum_{r=1}^m
w_r
\frac{a_rb_r(a_r-b_r)}
{(T_r+a_r+b_r)(T_r+a_r)(T_r+b_r)},
\end{equation}
where \(\Delta\) is the cost of \((a,b)\) minus the cost of \((b,a)\).
\end{lemma}

\begin{proof}
The network objective is a nonnegative weighted sum of single-resource
objectives. The algebra in \cref{lem:swap} remains valid when \(a_r=0\) or
\(b_r=0\) because \(T_r>0\). Summing the weighted resourcewise differences
gives \eqref{eq:netswap}.
\end{proof}

\begin{theorem}[Componentwise-chain ordering]
\label{thm:chain}
Suppose the item load vectors form a chain under componentwise order. Then any
componentwise nondecreasing order minimizes \(C^{\mathrm{net}}\), and the reverse
order maximizes it. In particular, if \(x_i=s_i v\) for a common profile
\(v\in\R_+^m\setminus\{0\}\) and scalars \(s_i>0\), sorting the scalar magnitudes
\(s_i\) increasingly is optimal.
\end{theorem}

\begin{proof}
If \(a\ge b\) componentwise, every summand in \eqref{eq:netswap} is nonnegative.
Swapping \((a,b)\) to \((b,a)\) therefore cannot increase cost. Repeated
adjacent interchanges prove the claim.
\end{proof}

For resources with \(q_r>0\), each resource has its own log-depletion vector.
Under the componentwise-chain assumption, the same adjacent swaps move every
such resource's log vector in the same majorization direction. Decreasing
componentwise order therefore minimizes every separable convex risk
\[
\sum_{r=1}^m w_r\sum_{k=1}^n g_r(d_{k,r})
\]
with convex \(g_r\), while increasing order minimizes the raw-fraction objective
\eqref{eq:networkobjective}.

The chain assumption is essentially the limit of any context-free sorting rule.

\begin{theorem}[Context reversal for pairwise priorities]
\label{thm:reversal}
For heterogeneous load vectors, the locally preferred order of a fixed pair can
reverse as the residual network state changes. Consequently, no priority
relation depending only on the two item vectors can agree with the optimal
adjacent order for all tail states.
\end{theorem}

\begin{proof}
Take two resources with equal weights and
\[
a=(4,1),\qquad b=(1,3).
\]
For tail state \(T=(1,1)\), \eqref{eq:netswap} gives
\[
\Delta(a,b;(1,1))
=
\frac{4\cdot1\cdot3}{6\cdot5\cdot2}
+
\frac{1\cdot3\cdot(-2)}{5\cdot2\cdot4}
=
\frac1{20}>0.
\]
Hence \((b,a)\) is locally cheaper. For tail state \(T=(10,1)\),
\[
\Delta(a,b;(10,1))
=
\frac{2}{385}-\frac3{20}
=-\frac{223}{1540}<0,
\]
so \((a,b)\) is locally cheaper. Since the same pair requires opposite orders in
two residual states, no context-free relation on item vectors alone can be
universally correct.
\end{proof}

This does not rule out state-dependent policies, dynamic programming,
approximation algorithms, or static indices on restricted load geometries. It
shows that no scalar priority can reproduce the correct adjacent comparison in
every residual state of the unrestricted network model.

\section{Economic and computing interpretations}
\label{sec:applications}

The model applies only when stage cost is genuinely proportional to the share of
capacity present immediately before the action. It is not a substitute for
travel cost, fixed setup cost, deadline penalties, price impact, or other
domain-specific effects.

\subsection{Computing systems}
Natural targets include staged retirement of servers, shards, replicas, or
service capacity; memory reclamation; and dependency-constrained
decommissioning. The raw-fraction objective favors removing small components
first. Convex functions of log depletion favor large-first schedules that smooth
multiplicative shocks. In a multi-resource system, loads may represent CPU,
memory, bandwidth, storage, or regional redundancy. \Cref{thm:reversal} says
heterogeneous systems generally need state-aware scheduling.

\subsection{Logistics and infrastructure}
Loads may represent warehouse zones, generation units, supplier capacity, or
temporary emergency resources. Precedence captures operational dependencies. The
stochastic theorem applies when uncertain loads sit in a likelihood-ratio-ordered
family, which is stronger than sorting means.

\subsection{Economics}
In a stylized market-exit or subsidy-withdrawal model, \(x_i/R_k\) is a
participant's share of the remaining system. Increasing order minimizes the sum
of ordinary percentage shocks. Decreasing order minimizes convex concentration
risk in continuously compounded shocks. A reserve \(q\) can represent outside
liquidity, a backstop, permanent capacity, or a diversified outside option.
Larger \(q\) weakens every pairwise ordering penalty in \eqref{eq:swap}.

\Cref{cor:phase} is the policy point: ``minimize disruption'' is incomplete until
disruption is specified as a concave aggregate of ordinary shares, an additive
log change, or a convex tail-risk measure.

\section{Limitations and open problems}
\label{sec:open}

\begin{enumerate}[leftmargin=2em]
\item \textbf{Complexity under precedence.} There is an exact ideal DP and a
greedy counterexample, but no NP-hardness classification or approximation
guarantee.
\item \textbf{Stochastic orders.} Likelihood-ratio order is strong.
Whether hazard-rate, dispersive, or usual stochastic order suffices for useful
subclasses is open.
\item \textbf{Network case.} Context reversal rules out a universal
context-free adjacent priority rule. Complexity and approximability of the
unrestricted multi-resource problem remain open.
\item \textbf{Empirical calibration.} Applications need evidence that fraction or
log shock is the actual cost, not just a convenient surrogate.
\item \textbf{Online models.} Items may arrive over time or reveal size only when
processed. Competitive analysis and regret bounds are natural next steps.
\end{enumerate}

\section{Conclusion}
\label{sec:conclusion}

A nested residual-fraction objective has an exact adjacent-interchange law:
small-first minimizes the raw cost, large-first maximizes it. The log
representation gives a stronger majorization fact: large-first smooths
multiplicative depletion, while small-first concentrates it. The same identity
drives an order-ideal DP under precedence, an MLR expectation theorem under
independent random sizes, and a componentwise sorting rule for multi-resource
chains. In heterogeneous networks, even a fixed pair can reverse local
preference with the residual state, marking a boundary for context-free
pairwise priorities.

Related cumulative residual-fraction sums already admit order-independent
integral bounds \citep{Pain2026}. The claim of this research is the permutation
theory for \eqref{eq:objective} and the package around one interchange identity,
not a new residual-ratio expression in isolation and not a new interchange
technique.

\section*{Declarations}

\subsection*{Funding}
This research received no external funding.

\subsection*{Competing interests}
The author declares no competing interests.

\subsection*{Data availability}
No datasets were generated or analysed in this study.

\subsection*{Code availability}
Companion Python and Kotlin verification scripts supplied with the manuscript
perform randomized numerical checks of the swap identity and majorization
theorem, enumerate the precedence counterexample, evaluate the two network
context states, and Monte Carlo check the stochastic theorem for exponential
sizes. The scripts are provided for transcription checks only and do not
replace the proofs.

\begingroup
\setlength{\bibsep}{0.5em}
\begin{thebibliography}{99}
\raggedright

\bibitem[Arnold(2007)]{Arnold2008}
Arnold~BC (2007) Majorization: Here, there and everywhere.
\emph{Statistical Science} 22(3):407--413.
\url{https://doi.org/10.1214/0883423060000000097}

\bibitem[Chang and Yao(1993)]{ChangYao1993}
Chang~CS, Yao~DD (1993) Rearrangement, majorization and stochastic scheduling.
\emph{Mathematics of Operations Research} 18(3):658--684.
\url{https://doi.org/10.1287/moor.18.3.658}

\bibitem[Correa and Schulz(2005)]{CorreaSchulz2005}
Correa~JR, Schulz~AS (2005) Single-machine scheduling with precedence
constraints.
\emph{Mathematics of Operations Research} 30(4):1005--1021.
\url{https://doi.org/10.1287/moor.1050.0158}

\bibitem[Hardy et~al.(1952)]{HardyLittlewoodPolya1952}
Hardy~GH, Littlewood~JE, P\'olya~G (1952) \emph{Inequalities}, 2nd~edn.
Cambridge University Press, Cambridge

\bibitem[J\"ager and Warode(2024)]{JagerWarode2024}
J\"ager~S, Warode~P (2024) Simple approximation algorithms for minimizing the
total weighted completion time of precedence-constrained jobs.
In: \emph{Proceedings of the 2024 Symposium on Simplicity in Algorithms (SOSA)}.
SIAM, pp~82--96.
\url{https://doi.org/10.1137/1.9781611977936.8}

\bibitem[Karlin and Rubin(1956)]{KarlinRubin1956}
Karlin~S, Rubin~H (1956) The theory of decision procedures for distributions
with monotone likelihood ratio.
\emph{The Annals of Mathematical Statistics} 27(2):272--299.
\url{https://doi.org/10.1214/aoms/1177728259}

\bibitem[Lawler(1978)]{Lawler1978}
Lawler~EL (1978) Sequencing jobs to minimize total weighted completion time
subject to precedence constraints.
\emph{Annals of Discrete Mathematics} 2:75--90.
\url{https://doi.org/10.1016/S0167-5060(08)70323-6}

\bibitem[Luce(1959)]{Luce1959}
Luce~RD (1959) \emph{Individual Choice Behavior: A Theoretical Analysis}.
Wiley, New York

\bibitem[Marshall et~al.(2011)]{MarshallOlkinArnold2011}
Marshall~AW, Olkin~I, Arnold~BC (2011) \emph{Inequalities: Theory of
Majorization and Its Applications}, 2nd~edn.
Springer, New York.
\url{https://doi.org/10.1007/978-0-387-68276-1}

\bibitem[Pain(2026)]{Pain2026}
Pain~JC (2026) Cumulative Riemann sums, distribution functions, and a universal
inequality. arXiv:2603.08959.
\url{https://doi.org/10.48550/arXiv.2603.08959}

\bibitem[Perman et~al.(1992)]{PermanPitmanYor1992}
Perman~M, Pitman~J, Yor~M (1992) Size-biased sampling of Poisson point processes
and excursions.
\emph{Probability Theory and Related Fields} 92(1):21--39.
\url{https://doi.org/10.1007/BF01205234}

\bibitem[Pitman and Tran(2015)]{PitmanTran2015}
Pitman~J, Tran~NM (2015) Size-biased permutation of a finite sequence with
independent and identically distributed terms.
\emph{Bernoulli} 21(4):2484--2512.
\url{https://doi.org/10.3150/14-BEJ652}

\bibitem[Plackett(1975)]{Plackett1975}
Plackett~RL (1975) The analysis of permutations.
\emph{Journal of the Royal Statistical Society: Series C (Applied Statistics)}
24(2):193--202.
\url{https://doi.org/10.2307/2346567}

\bibitem[Schaible and Shi(2003)]{SchaibleShi2003}
Schaible~S, Shi~J (2003) Fractional programming: The sum-of-ratios case.
\emph{Optimization Methods and Software} 18(2):219--229.
\url{https://doi.org/10.1080/1055678031000105242}

\bibitem[Shaked and Shanthikumar(2007)]{ShakedShanthikumar2007}
Shaked~M, Shanthikumar~JG (2007) \emph{Stochastic Orders}.
Springer, New York.
\url{https://doi.org/10.1007/978-0-387-34675-5}

\bibitem[Shanthikumar and Yao(1991)]{ShanthikumarYao1991}
Shanthikumar~JG, Yao~DD (1991) Bivariate characterization of some stochastic
order relations.
\emph{Advances in Applied Probability} 23(3):642--659.
\url{https://doi.org/10.2307/1427627}

\bibitem[Sidney(1975)]{Sidney1975}
Sidney~JB (1975) Decomposition algorithms for single-machine sequencing with
precedence relations and deferral costs.
\emph{Operations Research} 23(2):283--298.
\url{https://doi.org/10.1287/opre.23.2.283}

\bibitem[Smith(1956)]{Smith1956}
Smith~WE (1956) Various optimizers for single-stage production.
\emph{Naval Research Logistics Quarterly} 3(1--2):59--66.
\url{https://doi.org/10.1002/nav.3800030106}

\bibitem[Steiner(1984)]{Steiner1984}
Steiner~G (1984) Single machine scheduling with precedence constraints of
dimension 2.
\emph{Mathematics of Operations Research} 9(2):248--259.
\url{https://doi.org/10.1287/moor.9.2.248}

\end{thebibliography}
\endgroup

\end{document}
